\chapter{Peptic Ulcer Symptoms}

\section*{Definition}
Peptic ulcer disease (PUD) is a defect in the gastric or duodenal mucosa that extends through the muscularis mucosae, typically 5 mm or larger. It is classified as gastric or duodenal, with Helicobacter pylori infection and NSAID use accounting for the majority of cases.

\section*{What Happens in Your Body}
PUD results from an imbalance between aggressive factors (gastric acid, pepsin, bile salts, H. pylori urease, NSAIDs) and mucosal defense mechanisms (mucus-bicarbonate barrier, prostaglandins, epithelial renewal, blood flow). H. pylori colonizes the antrum, leading to hypergastrinemia and increased acid output.

\section*{Causes}
\begin{itemize}
  \item \textit{Helicobacter pylori} infection: 70--90\% of duodenal ulcers, 50--70\% of gastric ulcers
  \item \textbf{NSAIDs:} Non-selective (ibuprofen, naproxen, aspirin) and COX-2 selective inhibit chemical messengers that cause inflammation and pain synthesis, reducing mucosal blood flow and bicarbonate secretion
  \item \textbf{Pathologic hypersecretion:} Zollinger-Ellison syndrome (gastrinoma), systemic mastocytosis, short bowel syndrome
  \item \textbf{Other:} Smoking (delayed healing, increased recurrence), heavy alcohol use, critical illness (stress ulcers), radiation, chemotherapy (sunitinib, sorafenib), Crohn disease, CMV (immunocompromised)
  \item \textbf{Idiopathic:} Increasing prevalence in the era of declining H. pylori prevalence
\end{itemize}

\section*{When to See a Doctor}
See your doctor if:
\begin{itemize}
  \item Epigastric pain with hematemesis, coffee-ground emesis, melena, or hematochezia
  \item Sudden severe epigastric pain with peritonitis (suggests perforated ulcer)
  \item Persistent vomiting, dehydration, or weight loss
  \item Pain radiating to the back suggesting posterior penetrating ulcer (gastroduodenal artery erosion)
  \item Dysphagia, early satiety, or anemia without known blood loss
\end{itemize}

\section*{Related Symptoms}
\begin{itemize}
  \item Epigastric burning or gnawing pain (relieved by food in duodenal; worsened by food in gastric), bloating, nausea, early satiety, melena, hematemesis
\end{itemize}
\textbf{Important:} Classic symptom patterns (duodenal ulcer pain 2--5 hours after meals, gastric ulcer pain immediately after eating) are unreliable for distinguishing ulcer location and may be absent in up to 40\% of people.

\section*{Self-Care / Home Management}
\begin{itemize}
  \item Avoid NSAIDs, aspirin, alcohol, and smoking (correct potentially modifiable causes)
  \item Eat small, frequent meals; avoid known dietary triggers (spicy, acidic foods, caffeine)
  \item Antacids (calcium carbonate or aluminum/magnesium hydroxide) for short-term symptom relief
  \item If H. pylori testing is positive, complete the full course of prescribed eradication therapy
  \item Seek medical evaluation if symptoms persist beyond 2 weeks of empiric therapy
\end{itemize}

\section*{How Your Doctor Will Evaluate You}
\begin{itemize}
  \item EGD with biopsy is the gold standard for diagnosis; allows visualization, H. pylori testing (rapid urease test, histology), and biopsy to exclude malignancy
  \item H. pylori testing: urea breath test, stool antigen test, or biopsy-based testing; serology less reliable
  \item Barium upper GI series (double-contrast) if EGD unavailable (sensitivity 70--80\%)
  \item Fasting serum gastrin level if Zollinger-Ellison syndrome suspected (gastrin greater than 1000 pg/mL or elevated with secretin stimulation)
  \item Biopsy of all gastric ulcers (minimum 6--8 biopsies) to exclude adenocarcinoma; duodenal ulcers can be assumed benign
\end{itemize}

\section*{Treatment Options}
\begin{itemize}
  \item PPI (omeprazole 20 mg, pantoprazole 40 mg, esomeprazole 40 mg) for 4--8 weeks (duodenal) or 8--12 weeks (gastric)
  \item H. pylori eradication with 14-day therapy: triple (PPI, amoxicillin, clarithromycin) or bismuth-based quadruple therapy
  \item Discontinue NSAIDs; if must continue, co-prescribe PPI or misoprostol (200 mcg QID)
  \item Surgical management for perforation (omentopexy/Graham patch), bleeding (over-stitch with vessel ligation), or obstruction (pyloroplasty)
  \item Second-look endoscopy for high-risk bleeding ulcers (Forrest Ia--IIb) after endoscopic hemostasis (epinephrine injection, heater probe, clip placement)
\end{itemize}

\clearpage
